% Generated draft: requires mathematical and source review.
\documentclass[11pt]{article}
\usepackage[margin=25mm]{geometry}
\usepackage{fontspec}
\setmainfont{Latin Modern Roman}
\usepackage{amsmath,amssymb,mathtools,unicode-math}
\setmathfont{Latin Modern Math}
\usepackage{enumitem,needspace,xurl,hyperref}
\setlength{\parindent}{0pt}
\setlength{\parskip}{5pt}
\setlength{\emergencystretch}{3em}
\newcommand{\N}{\mathbb{N}}
\newcommand{\Z}{\mathbb{Z}}
\newcommand{\Q}{\mathbb{Q}}
\newcommand{\R}{\mathbb{R}}
\newcommand{\C}{\mathbb{C}}
\newcommand{\F}{\mathbb{F}}
\newcommand{\A}{\mathbb{A}}
\newcommand{\PP}{\mathbb{P}}
\newcommand{\E}{\mathbb{E}}
\newcommand{\eps}{\varepsilon}
\newcommand{\dd}{\,\mathrm{d}}
\begin{document}
\section*{UnsolvedMath ID: 30002450}
\textbf{Source problem code:} catalogue entry
\subsection{Definitions and mathematical statement}
The supplied question concerns quivers associated with triangulations of closed surfaces of positive genus having exactly two punctures. Under the conventional interpretation, let \(\Sigma\) be a compact connected oriented surface without boundary, of genus \(g\geq 1\), let \(M\subset\Sigma\) consist of exactly two distinct marked points, and let \(\tau\) be a triangulation of \((\Sigma,M)\). Write \(Q(\tau)\) for its signed adjacency quiver. The record does not explicitly specify orientability, connectedness, the kind of triangulation, or the conventions for self-folded triangles; this interpretation therefore requires source confirmation.

Fix a coefficient field \(k\), whose required hypotheses are not supplied. For a finite quiver \(Q\) without loops or oriented two-cycles, let \(k\langle\!\langle Q\rangle\!\rangle\) denote its completed path algebra, completed with respect to the ideal generated by the arrows. A potential on \(Q\) is a possibly infinite \(k\)-linear combination of cyclic paths, considered up to cyclic equivalence. Here cyclic equivalence identifies a cyclic path with its cyclic rotations and takes the closure of the resulting relations in the arrow-ideal topology.

Two potentials \(S,S'\) on \(Q\) are \emph{weakly right equivalent}, written \(S\sim_{\mathrm{wr}}S'\), if there are a scalar \(\lambda\in k^{\times}\) and a continuous \(k\)-algebra automorphism \(\varphi\) of \(k\langle\!\langle Q\rangle\!\rangle\), fixing every vertex idempotent, such that \(\varphi(S)\) is cyclically equivalent to \(\lambda S'\). A potential is \emph{non-degenerate} in the quiver-with-potential mutation sense if every finite sequence of vertex mutations can be performed successively and, after reduction at each step, the resulting quiver has no oriented two-cycles. Let \(\mathcal P_{\mathrm{nd}}(Q;k)\) denote the collection of such potentials.

\textbf{Question, with unresolved source conventions.} For every surface and triangulation in the intended class, is there exactly one weak right equivalence class of non-degenerate potentials on its associated quiver? With the coefficient field and triangulation conventions fixed, the question is
\[
\left|\mathcal P_{\mathrm{nd}}\bigl(Q(\tau);k\bigr)/\!\sim_{\mathrm{wr}}\right|=1
\qquad\text{for every admissible }(\Sigma,M,\tau),\quad g(\Sigma)\geq 1,\quad |M|=2.
\]
This asks both for existence and for equivalence of any two non-degenerate potentials on each fixed quiver. It does not assert equivalence between potentials on different quivers. The coefficient-field hypotheses and the precise class of admissible triangulations are unresolved in the supplied record; the display is a formulation schema, not an assertion universally quantified over all fields.
\subsection{Short English statement}
For every quiver obtained from a triangulation of a closed positive-genus surface with exactly two punctures, ask whether a non-degenerate potential exists and whether any two such potentials differ only by a right equivalence and multiplication by a nonzero scalar. The coefficient field and triangulation conventions require confirmation. The supplied record labels the problem open; this draft does not independently verify that status.
\subsection{Sources}
\begin{itemize}
\item Supplied local UnsolvedMath registry record, permanent ID \(30002450\), title: \emph{Potentials for Quivers from Twice-Punctured Surfaces}; collection: \emph{Oberwolfach Reports Open Problems}; exact source problem code: \texttt{catalogue entry}. The record gives proposed year \(2013\) and names no proposer.
\item Supplied source citation: \emph{Cluster Algebras and Related Topics (2014)}, Oberwolfach Reports, pp.~3379--3432, DOI \texttt{10.4171/owr/2013/58}. \url{https://doi.org/10.4171/owr/2013/58}. This is also the supplied literature-source link labeled ``original OWR report.'' The background describes Oberwolfach Reports as documenting mathematical workshops and their open-problem discussions.
\item Supplied literature-source attribution: Geiss, Labardini-Fragoso and Schroer, DOI \texttt{10.1016/j.aim.2015.09.038}. \url{https://doi.org/10.1016/j.aim.2015.09.038}. No article title or precise theorem reference is supplied.
\item Literature assessment dated exactly \texttt{2026-08-21}: the supplied curation labels the problem open and reports that uniqueness is proved for many mutation classes, while closed positive-genus surfaces with exactly two punctures remain an unresolved classification case. This is supplied dated curation derived from the OWR material, not an independently verified assessment; the status and cited claims require independent verification.
\item The registry labels the extracted statement ``exact'' and reports that manual review found it self-contained, mathematically coherent, and free of LaTeX or report-layout artifacts. These are supplied verification claims, not verification performed for this draft.
\end{itemize}
\end{document}
